%% ma: L12-B14-0066
%% lop: 12
%% chuong: 5
%% bai: 14
%% dang: 12.14.21
%% loai: TN4
%% muc: NB
%% dap_an: "D"
%% nguon: "To_n_12_PT_M_t_Ph_ng.pdf (Đề kiểm tra đánh giá kiến thức, Chuyên đề Phương trình mặt phẳng, Nguyễn Anh Khôi - Lê Minh Khôi), Câu 21"
%% kien_thuc: "Bài 14 (phương trình mặt phẳng theo đoạn chắn)"
%% trang_thai: cho-duyet
%% ly_do: "Dạng 12.14.21 mới đề xuất, chờ giáo viên duyệt."
%% ghi_chu: "Tài liệu gốc không kèm đáp án; đáp án do Claude giải và kiểm bằng sympy."
\begin{ex}
Trong không gian $Oxyz$, cho ba điểm $A(2;0;0)$, $B(0;-1;0)$, $C(0;0;3)$. Mặt phẳng $(ABC)$ có phương trình là
\choice{$\dfrac{x}{-2}+\dfrac{y}{1}+\dfrac{z}{3}=1$}{$\dfrac{x}{2}+\dfrac{y}{1}+\dfrac{z}{-3}=1$}{$\dfrac{x}{2}+\dfrac{y}{1}+\dfrac{z}{3}=1$}{\True $\dfrac{x}{2}+\dfrac{y}{-1}+\dfrac{z}{3}=1$}
\loigiai{Phương trình theo đoạn chắn: $(ABC)\colon \dfrac{x}{2}+\dfrac{y}{-1}+\dfrac{z}{3}=1$. Chọn D.}
\end{ex}
